% META: {"id": "1NSI-ADGK-CO-020", "chapitre": "1NSI-ALGO-PARCOURS-TRIS", "type_objet": "correction", "exercice_ref": "1NSI-ADGK-EX-020", "capacites_codes": ["C4"], "sources_inspiration": [], "status": "approved", "capacites": ["1NSI-ALGO-PARCOURS-TRIS-C4"], "parcours": 2, "duree_min": 2, "fichier_exercice": "chapitres/1NSI-ALGO-PARCOURS-TRIS/exercices/1NSI-ADGK-EX-020.tex", "mode_creation": "ex_nihilo"}
% BEGIN-VERIFY
% from sympy import *
% x = symbols('x')
% g = exp(x) - 1 - x
% g2 = diff(g, x, 2)
% assert g2 == exp(x)
% # g2 = exp(x) > 0 pour tout x (convexe)
% assert g.subs(x, 0) == 0
% assert diff(g, x).subs(x, 0) == 0
% # donc g >= 0 car convexe et minimum en 0
% assert g.subs(x, 1) == exp(1) - 2
% assert g.subs(x, -1) == exp(-1)
% END-VERIFY
\begin{corrige}{1NSI-ADGK-EX-020}

\commentaireMarge{Etudier $g(x) = \mathrm{e}^x - 1 - x$ et utiliser la convexite.}

Posons $g(x) = \mathrm{e}^x - 1 - x$. On a $g'(x) = \mathrm{e}^x - 1$ et $g''(x) = \mathrm{e}^x > 0$.

La fonction $g$ est donc \textbf{convexe} sur $\mathbb{R}$. Son minimum est atteint en $x=0$ (car $g'(0)=0$), et $g(0) = 0$.

Par convexite, $g(x) \geq g(0) = 0$ pour tout $x$, soit $\mathrm{e}^x \geq 1 + x$.

\end{corrige}
